\documentclass[10pt]{article}
\usepackage[margin=0.9in]{geometry}
\usepackage{enumitem}
\usepackage{hyperref}
\setlist{nosep,leftmargin=*}
\setlength{\parskip}{4pt}
\setlength{\parindent}{0pt}

\title{\vspace{-1.5em}Matrix-Multiply Accelerator Integration in a RISC-V SoC\vspace{-0.5em}}
\author{Progress Report (continued)}
\date{\vspace{-1em}\today}

\begin{document}
\maketitle
\vspace{-1em}

\section*{Accelerator Design}
The first version of the accelerator computed products on a fixed $4\times4$ array and
moved data over the bus one word at a time. Because each transaction waited for its reply
before the next began, it could only handle matrices of at most $4\times4$, and most of
the time was spent waiting for memory. The design was then reworked in three ways.

The array size is now a build-time parameter, along with the depth of the input buffers.
The same design can be built with arrays of different sizes, including sizes that are not
powers of two, and the software reads the hardware configuration at run time, so one test
program runs unchanged on every build.

Moving one word per transaction was the main bottleneck, so the read and write paths were
rebuilt to use AXI bursts, which transfer a whole row of data in one request. Several reads
can be in flight at once, and bursts are split so that none crosses a 4\,KB boundary, as the
bus protocol requires. This brought the accelerator much closer to moving one word per
cycle.

Hardware tiling then lets a single start command compute a product of any size. The result
is divided into tiles, and the inner dimension is processed in chunks whose partial sums stay
in the array, so they never go back to memory. Matrices can also be read as sub-matrices of
larger ones through a row stride. When the inner dimension fits in one chunk, the accelerator
reuses each row of $A$ across all the tiles that need it, which reduces the memory traffic of
larger products.

\section*{Hardware Testing and Verification}
Verification follows the same approach at each stage. The processor computes every product
with an ordinary triple loop over the same input matrices, and the accelerator computes it in
hardware. The test program compares the two results element by element. The test covers a range of matrix shapes, edge cases and invalid inputs, and it was repeated
across several build settings.

The $4\times4$ version was verified first, and the tiled version was then verified in the same
way, so that the tiling logic could be checked against the simpler design before any larger
products were attempted. Both were checked in simulation before any hardware work.

The design was then run on a RealDigital Boolean board, which carries a Xilinx XC7S50 FPGA. The
board build used a $5\times5$ array. The board runs the same test program, and its output is
cross-checked against the processor's own result on the board. The board output also matches a
Verilator simulation of the same design, character for character. The design meets timing at
the board's 25\,MHz clock, and it runs the test suite with products up to $64\times64$, with and
without reuse.

To improve performance, the maximum K chunk size (\texttt{MAX\_K}) was varied. A larger chunk
reduces the number of chunks each tile must go through, and it increases the reuse of loaded
data, so the cycle counts drop noticeably. The cost is the buffer size: each increase in
\texttt{MAX\_K} adds flip-flops and multiplexer logic. On the board, raising the value from the
initial setting to 64 reduced the cycle count of the larger tests while keeping timing and
device usage within limits.

\section*{Remaining Work and Limitations}
The main limitation is memory. The Boolean board has no external RAM, so its main memory is the
on-chip block RAM, about 128\,KB. This holds the program and all matrices, so the largest products
that fit are in the tens, and products of $256\times256$ and above cannot run on this board. The
memory decodes only its low address bits, so an address beyond its end silently wraps around and
can overwrite the program. The test does not yet check that its buffers fit. Other limits include
the bandwidth of the single shared bus, which keeps the accelerator memory-bound, the growth of
the input buffers in logic as \texttt{MAX\_K} increases, and the 16-bit dimension registers.
Completion is detected by polling rather than an interrupt, and a bus error during a job is not
handled cleanly.

The next steps are to run products of $256\times256$ and larger on the Arty A7 board, which has
external DDR3 memory and an existing memory project in the repository. On the Boolean board, a
bounds check should be added before any larger test. A software baseline timed with the
processor's cycle counter is needed for a fair speed-up comparison, and the robustness items
above should be completed.

\end{document}
